% !PRINTED 460
\weberpage{460}
%<<CONT
which, developed, takes the form
\weq{2}{\Delta = 5x_1^2x_2^2x_3^2 - x_1x_3^5 - x_2x_1^5 - x_3x_2^5.}
On the curve $\Delta$, then, lie the vertices of the coordinate triangle,
and consequently all the sevenfold poles $P_7$ lie on $\Delta$ and
form the complete system of the points of intersection of $f$ and $\Delta$.

A further covariant, namely one of degree 14,
we obtain from \S.\,87, (3), namely the determinant:
\weq{3}{C = \frac{1}{9}
\begin{vmatrix}
f''(x_1,x_1), & f''(x_1,x_2), & f''(x_1,x_3), & \Delta'(x_1) \\
f''(x_2,x_1), & f''(x_2,x_2), & f''(x_2,x_3), & \Delta'(x_2) \\
f''(x_3,x_1), & f''(x_3,x_2), & f''(x_3,x_3), & \Delta'(x_3) \\
\Delta'(x_1), & \Delta'(x_2), & \Delta'(x_3), & 0
\end{vmatrix}.}
The determinant $C$ can, though somewhat laboriously, be computed
from (3), and takes the expression:
\weq{4}{C = \Sigma x_1^{14} - 34x_1x_2x_3\,\Sigma x_1^{10}x_3 - 250x_1x_2x_3\,\Sigma x_1^8x_2^3 + 375x_1^2x_2^2x_3^2\,\Sigma x_1^6x_2^2 + 18\,\Sigma x_1^7x_2^7 - 126x_1^3x_2^3x_3^3\,\Sigma x_1^3x_2^2,}
in which the symbol $\Sigma$ means that the sum is to be taken of the
three terms which are obtained from the first when $x_1, x_2, x_3$
are cyclically interchanged. The computation of the first term
$x_1^{14}$ already suffices to show that the curve $C = 0$
does \gesperrt{not} pass through the vertices of the coordinate triangle.

A fourth invariant, namely one of degree 21,
we obtain according to \S.\,87, (6), by forming the functional determinant
of the three forms $f$, $\Delta$, $C$:
\weq{5}{K = \frac{1}{14}
\begin{vmatrix}
f'(x_1), & \Delta'(x_1), & C'(x_1) \\
f'(x_2), & \Delta'(x_2), & C'(x_2) \\
f'(x_3), & \Delta'(x_3), & C'(x_3)
\end{vmatrix},}
which likewise can be computed from this. We give here only the
first three terms, from which one sees that this curve too
does not pass through the vertices of the coordinate triangle
\weq{6}{K = x_1^{21} + x_2^{21} + x_3^{21} + \cdots{}^{1)}.}

$^{1)}$ The fully computed expressions are to be found in the
memoir of \gesperrt{Gordan}: ``Ueber die typische Darstellung
der ternären biquadratischen Form $f = x_1^3x_2 + x_2^3x_3 + x_3^3x_1$''
[Mathem. Annalen, Bd. XVII, S. 366 (1880)]. Also in \gesperrt{Klein-Fricke},
Modulfunctionen, Bd. I, S. 734.
